% =============================================================================
% Ethics Statement and Broader Impact
% Self-contained file: % =============================================================================
% Ethics Statement and Broader Impact
% Self-contained file: % =============================================================================
% Ethics Statement and Broader Impact
% Self-contained file: % =============================================================================
% Ethics Statement and Broader Impact
% Self-contained file: \input{ethics_statement} into main.tex
% Intended placement: after the Conclusion section, before References.
% Label: \label{sec:impact}
% =============================================================================

\section{Broader Impact and Ethics Statement}
\label{sec:impact}

\subsection{Broader Impact}

\paragraph{Positive impacts.}
\ours{} lowers the barrier to reproducible RL post-training research by
providing a unified benchmark harness that works across multiple compute
platforms.  Prior work frequently reports results on a single proprietary
GPU cluster, making it difficult to determine how much of the observed
performance is attributable to the training algorithm versus the hardware
and scheduling environment.  By running identical experiments on both a
closed API (Tinker) and an open, explicitly provisioned cloud (Modal), we
surface \emph{platform-dependent variance} as a first-class experimental
variable.  This transparency allows practitioners to make more informed
decisions when selecting infrastructure for RL fine-tuning and helps the
community calibrate trust in benchmark results from single-platform papers.

Releasing all code, model checkpoints, and experiment logs publicly enables
independent replication, third-party auditing, and incremental improvement
without the need to re-run expensive GPU-hours from scratch.  We expect the
benchmark to be most valuable for researchers at resource-constrained
institutions who need a principled starting point for RL post-training
experiments.

\paragraph{Potential negative impacts and limitations.}
\begin{itemize}
  \item \textbf{Reward hacking.}  RLHF and GRPO optimise a proxy reward
        signal rather than the true underlying objective.  Models trained
        to maximise verifier-based rewards on GSM8K or HumanEval may learn
        to exploit idiosyncrasies of the verification function rather than
        genuinely improving reasoning or coding ability.  We include KL
        divergence tracking precisely to monitor policy drift, but reward
        hacking remains a risk in any production deployment of RL-fine-tuned
        models.

  \item \textbf{Alignment concerns.}  Reinforcement learning from human or
        proxy feedback can shift model behaviour in directions that are
        difficult to predict or reverse.  Models trained with poorly
        specified reward signals may become better at appearing correct while
        simultaneously becoming less calibrated or more prone to confident
        errors.  Our benchmark focuses on mathematically verifiable tasks,
        which partially mitigates this risk, but we caution against
        uncritical transfer of RL post-training pipelines to open-ended or
        subjective task domains without additional alignment safeguards.

  \item \textbf{Compute access disparities.}  Even with cost-optimised
        LoRA fine-tuning, reproducing the full \ours{} benchmark requires
        access to NVIDIA H100 GPUs (or equivalent) for 200--600 GPU-hours.
        This remains out of reach for many research groups, particularly in
        the Global South.  We partially address this by providing pretrained
        checkpoints, but the barrier to running novel experiments remains
        high.

  \item \textbf{Closed-platform dependency.}  A portion of our benchmark
        relies on the Tinker API, a proprietary, closed-source service.
        Results from Tinker experiments cannot be independently reproduced
        by researchers without API access.  We consider this a significant
        limitation and flag all Tinker-derived results accordingly in our
        figures.  We strongly recommend that readers treat Modal/open-source
        results as the primary reproducibility reference.
\end{itemize}

% ---------------------------------------------------------------------------
\subsection{Compute Costs and Carbon Footprint}
% ---------------------------------------------------------------------------

Table~\ref{tab:carbon} summarises our compute allocation and estimated
carbon footprint.  Carbon intensity estimates follow the methodology of
\citet{patterson2024empirical}, using a conservative average grid intensity
of 0.233 kg\,CO$_2$-eq per kWh for the United States (EPA eGRID 2023
average) and a power usage effectiveness (PUE) of 1.1 for cloud data
centres.  NVIDIA H100 SXM5 TDP is 700\,W; NVIDIA L4 TDP is 72\,W.

\begin{table}[h]
\centering
\caption{Compute allocation and estimated carbon footprint by platform.
  Tinker API carbon figures are marked N/A because GPU type is undisclosed.
  Modal figures use confirmed hardware specifications.}
\label{tab:carbon}
\begin{tabular}{llcccc}
\toprule
\textbf{Platform} & \textbf{GPU} & \textbf{Est. GPU-h} &
\textbf{TDP (W)} & \textbf{Est. kWh} &
\textbf{Est. CO\textsubscript{2} (kg)} \\
\midrule
Modal (H100 SXM5) & H100 SXM5 & $\sim$250 & 700 & $\sim$192 & $\sim$0.49 \\
Modal (L4, older runs) & L4       & $\sim$50  & 72  & $\sim$4   & $\sim$0.01 \\
Tinker API            & unknown  & $\sim$950 & N/A & N/A       & N/A        \\
\midrule
\multicolumn{3}{l}{\textbf{Total (Modal only, estimable)}}      & & $\sim$196 & $\sim$0.50 \\
\bottomrule
\end{tabular}
\end{table}

The estimated carbon footprint for the Modal portion of our experiments is
approximately \textbf{0.4--0.6\,kg\,CO$_2$-eq}, comparable to a short
car journey.  We cannot estimate the Tinker API footprint due to the
platform's opacity.  We acknowledge that the full experimental campaign
(including failed runs) consumed additional GPU-hours; our estimates above
reflect \emph{successful} jobs only.

% ---------------------------------------------------------------------------
\subsection{Data Provenance}
% ---------------------------------------------------------------------------

All training and evaluation datasets used in \ours{} are publicly available
research datasets.  No private, proprietary, or sensitive data were used
at any stage.

\begin{itemize}
  \item \textbf{GSM8K}~\cite{cobbe2021training}: 8,500 grade-school math
        word problems with natural-language solutions, released by OpenAI
        under the MIT licence.  We use the standard train/test splits
        without modification.

  \item \textbf{HumanEval}: 164 Python programming problems with unit
        tests, released by OpenAI under the MIT licence.  We use the
        original evaluation protocol (pass@$k$) without modification.

  \item \textbf{Synthetic tool-use dataset}: Generated by the authors from
        publicly documented APIs (calculator, web-search stub, calendar).
        No third-party proprietary data or user data were used in
        generation.  The generation prompts and all resulting examples are
        included in the repository under MIT licence.
\end{itemize}

No web-scraped data, personally identifiable information, or data subject
to copyright restrictions were used.  No human annotators were employed.

% ---------------------------------------------------------------------------
\subsection{Limitation Acknowledgment: Tinker API Closed-Source Nature}
% ---------------------------------------------------------------------------

A central tension in \ours{} is that we benchmark \emph{against} a
closed-source platform while simultaneously advocating for reproducibility.
We include Tinker API experiments for three reasons: (i) Tinker represents
a realistic production environment for practitioners who lack on-premises
GPU clusters; (ii) the platform/open-source gap is itself a scientifically
interesting quantity to measure; and (iii) the high failure rate (11 of 14
experiments, 79\%) is a finding worth documenting publicly, as it reveals
the fragility of API-dependent training pipelines.

We make the following commitments to mitigate the reproducibility gap:
\begin{enumerate}
  \item All Tinker API experiment scripts, configuration files, and raw
        logs are archived in the repository.  Researchers with Tinker
        access can attempt replication using our exact configurations.
  \item All figures derived \emph{solely} from Tinker data are marked with
        ``\dag{}'' and accompanied by a note that independent reproduction
        requires Tinker API access.
  \item Primary statistical conclusions are drawn from Modal experiments
        wherever possible.  Tinker-only results are reported descriptively
        and are not used to support quantitative claims.
  \item If Tinker releases an open-source version or a compatible open
        alternative becomes available, we commit to re-running failed
        experiments and issuing a revised version of this paper.
\end{enumerate}

We believe honest disclosure of these constraints is more valuable to the
community than suppressing the Tinker experiments or presenting only
successful runs.
 into main.tex
% Intended placement: after the Conclusion section, before References.
% Label: \label{sec:impact}
% =============================================================================

\section{Broader Impact and Ethics Statement}
\label{sec:impact}

\subsection{Broader Impact}

\paragraph{Positive impacts.}
\ours{} lowers the barrier to reproducible RL post-training research by
providing a unified benchmark harness that works across multiple compute
platforms.  Prior work frequently reports results on a single proprietary
GPU cluster, making it difficult to determine how much of the observed
performance is attributable to the training algorithm versus the hardware
and scheduling environment.  By running identical experiments on both a
closed API (Tinker) and an open, explicitly provisioned cloud (Modal), we
surface \emph{platform-dependent variance} as a first-class experimental
variable.  This transparency allows practitioners to make more informed
decisions when selecting infrastructure for RL fine-tuning and helps the
community calibrate trust in benchmark results from single-platform papers.

Releasing all code, model checkpoints, and experiment logs publicly enables
independent replication, third-party auditing, and incremental improvement
without the need to re-run expensive GPU-hours from scratch.  We expect the
benchmark to be most valuable for researchers at resource-constrained
institutions who need a principled starting point for RL post-training
experiments.

\paragraph{Potential negative impacts and limitations.}
\begin{itemize}
  \item \textbf{Reward hacking.}  RLHF and GRPO optimise a proxy reward
        signal rather than the true underlying objective.  Models trained
        to maximise verifier-based rewards on GSM8K or HumanEval may learn
        to exploit idiosyncrasies of the verification function rather than
        genuinely improving reasoning or coding ability.  We include KL
        divergence tracking precisely to monitor policy drift, but reward
        hacking remains a risk in any production deployment of RL-fine-tuned
        models.

  \item \textbf{Alignment concerns.}  Reinforcement learning from human or
        proxy feedback can shift model behaviour in directions that are
        difficult to predict or reverse.  Models trained with poorly
        specified reward signals may become better at appearing correct while
        simultaneously becoming less calibrated or more prone to confident
        errors.  Our benchmark focuses on mathematically verifiable tasks,
        which partially mitigates this risk, but we caution against
        uncritical transfer of RL post-training pipelines to open-ended or
        subjective task domains without additional alignment safeguards.

  \item \textbf{Compute access disparities.}  Even with cost-optimised
        LoRA fine-tuning, reproducing the full \ours{} benchmark requires
        access to NVIDIA H100 GPUs (or equivalent) for 200--600 GPU-hours.
        This remains out of reach for many research groups, particularly in
        the Global South.  We partially address this by providing pretrained
        checkpoints, but the barrier to running novel experiments remains
        high.

  \item \textbf{Closed-platform dependency.}  A portion of our benchmark
        relies on the Tinker API, a proprietary, closed-source service.
        Results from Tinker experiments cannot be independently reproduced
        by researchers without API access.  We consider this a significant
        limitation and flag all Tinker-derived results accordingly in our
        figures.  We strongly recommend that readers treat Modal/open-source
        results as the primary reproducibility reference.
\end{itemize}

% ---------------------------------------------------------------------------
\subsection{Compute Costs and Carbon Footprint}
% ---------------------------------------------------------------------------

Table~\ref{tab:carbon} summarises our compute allocation and estimated
carbon footprint.  Carbon intensity estimates follow the methodology of
\citet{patterson2024empirical}, using a conservative average grid intensity
of 0.233 kg\,CO$_2$-eq per kWh for the United States (EPA eGRID 2023
average) and a power usage effectiveness (PUE) of 1.1 for cloud data
centres.  NVIDIA H100 SXM5 TDP is 700\,W; NVIDIA L4 TDP is 72\,W.

\begin{table}[h]
\centering
\caption{Compute allocation and estimated carbon footprint by platform.
  Tinker API carbon figures are marked N/A because GPU type is undisclosed.
  Modal figures use confirmed hardware specifications.}
\label{tab:carbon}
\begin{tabular}{llcccc}
\toprule
\textbf{Platform} & \textbf{GPU} & \textbf{Est. GPU-h} &
\textbf{TDP (W)} & \textbf{Est. kWh} &
\textbf{Est. CO\textsubscript{2} (kg)} \\
\midrule
Modal (H100 SXM5) & H100 SXM5 & $\sim$250 & 700 & $\sim$192 & $\sim$0.49 \\
Modal (L4, older runs) & L4       & $\sim$50  & 72  & $\sim$4   & $\sim$0.01 \\
Tinker API            & unknown  & $\sim$950 & N/A & N/A       & N/A        \\
\midrule
\multicolumn{3}{l}{\textbf{Total (Modal only, estimable)}}      & & $\sim$196 & $\sim$0.50 \\
\bottomrule
\end{tabular}
\end{table}

The estimated carbon footprint for the Modal portion of our experiments is
approximately \textbf{0.4--0.6\,kg\,CO$_2$-eq}, comparable to a short
car journey.  We cannot estimate the Tinker API footprint due to the
platform's opacity.  We acknowledge that the full experimental campaign
(including failed runs) consumed additional GPU-hours; our estimates above
reflect \emph{successful} jobs only.

% ---------------------------------------------------------------------------
\subsection{Data Provenance}
% ---------------------------------------------------------------------------

All training and evaluation datasets used in \ours{} are publicly available
research datasets.  No private, proprietary, or sensitive data were used
at any stage.

\begin{itemize}
  \item \textbf{GSM8K}~\cite{cobbe2021training}: 8,500 grade-school math
        word problems with natural-language solutions, released by OpenAI
        under the MIT licence.  We use the standard train/test splits
        without modification.

  \item \textbf{HumanEval}: 164 Python programming problems with unit
        tests, released by OpenAI under the MIT licence.  We use the
        original evaluation protocol (pass@$k$) without modification.

  \item \textbf{Synthetic tool-use dataset}: Generated by the authors from
        publicly documented APIs (calculator, web-search stub, calendar).
        No third-party proprietary data or user data were used in
        generation.  The generation prompts and all resulting examples are
        included in the repository under MIT licence.
\end{itemize}

No web-scraped data, personally identifiable information, or data subject
to copyright restrictions were used.  No human annotators were employed.

% ---------------------------------------------------------------------------
\subsection{Limitation Acknowledgment: Tinker API Closed-Source Nature}
% ---------------------------------------------------------------------------

A central tension in \ours{} is that we benchmark \emph{against} a
closed-source platform while simultaneously advocating for reproducibility.
We include Tinker API experiments for three reasons: (i) Tinker represents
a realistic production environment for practitioners who lack on-premises
GPU clusters; (ii) the platform/open-source gap is itself a scientifically
interesting quantity to measure; and (iii) the high failure rate (11 of 14
experiments, 79\%) is a finding worth documenting publicly, as it reveals
the fragility of API-dependent training pipelines.

We make the following commitments to mitigate the reproducibility gap:
\begin{enumerate}
  \item All Tinker API experiment scripts, configuration files, and raw
        logs are archived in the repository.  Researchers with Tinker
        access can attempt replication using our exact configurations.
  \item All figures derived \emph{solely} from Tinker data are marked with
        ``\dag{}'' and accompanied by a note that independent reproduction
        requires Tinker API access.
  \item Primary statistical conclusions are drawn from Modal experiments
        wherever possible.  Tinker-only results are reported descriptively
        and are not used to support quantitative claims.
  \item If Tinker releases an open-source version or a compatible open
        alternative becomes available, we commit to re-running failed
        experiments and issuing a revised version of this paper.
\end{enumerate}

We believe honest disclosure of these constraints is more valuable to the
community than suppressing the Tinker experiments or presenting only
successful runs.
 into main.tex
% Intended placement: after the Conclusion section, before References.
% Label: \label{sec:impact}
% =============================================================================

\section{Broader Impact and Ethics Statement}
\label{sec:impact}

\subsection{Broader Impact}

\paragraph{Positive impacts.}
\ours{} lowers the barrier to reproducible RL post-training research by
providing a unified benchmark harness that works across multiple compute
platforms.  Prior work frequently reports results on a single proprietary
GPU cluster, making it difficult to determine how much of the observed
performance is attributable to the training algorithm versus the hardware
and scheduling environment.  By running identical experiments on both a
closed API (Tinker) and an open, explicitly provisioned cloud (Modal), we
surface \emph{platform-dependent variance} as a first-class experimental
variable.  This transparency allows practitioners to make more informed
decisions when selecting infrastructure for RL fine-tuning and helps the
community calibrate trust in benchmark results from single-platform papers.

Releasing all code, model checkpoints, and experiment logs publicly enables
independent replication, third-party auditing, and incremental improvement
without the need to re-run expensive GPU-hours from scratch.  We expect the
benchmark to be most valuable for researchers at resource-constrained
institutions who need a principled starting point for RL post-training
experiments.

\paragraph{Potential negative impacts and limitations.}
\begin{itemize}
  \item \textbf{Reward hacking.}  RLHF and GRPO optimise a proxy reward
        signal rather than the true underlying objective.  Models trained
        to maximise verifier-based rewards on GSM8K or HumanEval may learn
        to exploit idiosyncrasies of the verification function rather than
        genuinely improving reasoning or coding ability.  We include KL
        divergence tracking precisely to monitor policy drift, but reward
        hacking remains a risk in any production deployment of RL-fine-tuned
        models.

  \item \textbf{Alignment concerns.}  Reinforcement learning from human or
        proxy feedback can shift model behaviour in directions that are
        difficult to predict or reverse.  Models trained with poorly
        specified reward signals may become better at appearing correct while
        simultaneously becoming less calibrated or more prone to confident
        errors.  Our benchmark focuses on mathematically verifiable tasks,
        which partially mitigates this risk, but we caution against
        uncritical transfer of RL post-training pipelines to open-ended or
        subjective task domains without additional alignment safeguards.

  \item \textbf{Compute access disparities.}  Even with cost-optimised
        LoRA fine-tuning, reproducing the full \ours{} benchmark requires
        access to NVIDIA H100 GPUs (or equivalent) for 200--600 GPU-hours.
        This remains out of reach for many research groups, particularly in
        the Global South.  We partially address this by providing pretrained
        checkpoints, but the barrier to running novel experiments remains
        high.

  \item \textbf{Closed-platform dependency.}  A portion of our benchmark
        relies on the Tinker API, a proprietary, closed-source service.
        Results from Tinker experiments cannot be independently reproduced
        by researchers without API access.  We consider this a significant
        limitation and flag all Tinker-derived results accordingly in our
        figures.  We strongly recommend that readers treat Modal/open-source
        results as the primary reproducibility reference.
\end{itemize}

% ---------------------------------------------------------------------------
\subsection{Compute Costs and Carbon Footprint}
% ---------------------------------------------------------------------------

Table~\ref{tab:carbon} summarises our compute allocation and estimated
carbon footprint.  Carbon intensity estimates follow the methodology of
\citet{patterson2024empirical}, using a conservative average grid intensity
of 0.233 kg\,CO$_2$-eq per kWh for the United States (EPA eGRID 2023
average) and a power usage effectiveness (PUE) of 1.1 for cloud data
centres.  NVIDIA H100 SXM5 TDP is 700\,W; NVIDIA L4 TDP is 72\,W.

\begin{table}[h]
\centering
\caption{Compute allocation and estimated carbon footprint by platform.
  Tinker API carbon figures are marked N/A because GPU type is undisclosed.
  Modal figures use confirmed hardware specifications.}
\label{tab:carbon}
\begin{tabular}{llcccc}
\toprule
\textbf{Platform} & \textbf{GPU} & \textbf{Est. GPU-h} &
\textbf{TDP (W)} & \textbf{Est. kWh} &
\textbf{Est. CO\textsubscript{2} (kg)} \\
\midrule
Modal (H100 SXM5) & H100 SXM5 & $\sim$250 & 700 & $\sim$192 & $\sim$0.49 \\
Modal (L4, older runs) & L4       & $\sim$50  & 72  & $\sim$4   & $\sim$0.01 \\
Tinker API            & unknown  & $\sim$950 & N/A & N/A       & N/A        \\
\midrule
\multicolumn{3}{l}{\textbf{Total (Modal only, estimable)}}      & & $\sim$196 & $\sim$0.50 \\
\bottomrule
\end{tabular}
\end{table}

The estimated carbon footprint for the Modal portion of our experiments is
approximately \textbf{0.4--0.6\,kg\,CO$_2$-eq}, comparable to a short
car journey.  We cannot estimate the Tinker API footprint due to the
platform's opacity.  We acknowledge that the full experimental campaign
(including failed runs) consumed additional GPU-hours; our estimates above
reflect \emph{successful} jobs only.

% ---------------------------------------------------------------------------
\subsection{Data Provenance}
% ---------------------------------------------------------------------------

All training and evaluation datasets used in \ours{} are publicly available
research datasets.  No private, proprietary, or sensitive data were used
at any stage.

\begin{itemize}
  \item \textbf{GSM8K}~\cite{cobbe2021training}: 8,500 grade-school math
        word problems with natural-language solutions, released by OpenAI
        under the MIT licence.  We use the standard train/test splits
        without modification.

  \item \textbf{HumanEval}: 164 Python programming problems with unit
        tests, released by OpenAI under the MIT licence.  We use the
        original evaluation protocol (pass@$k$) without modification.

  \item \textbf{Synthetic tool-use dataset}: Generated by the authors from
        publicly documented APIs (calculator, web-search stub, calendar).
        No third-party proprietary data or user data were used in
        generation.  The generation prompts and all resulting examples are
        included in the repository under MIT licence.
\end{itemize}

No web-scraped data, personally identifiable information, or data subject
to copyright restrictions were used.  No human annotators were employed.

% ---------------------------------------------------------------------------
\subsection{Limitation Acknowledgment: Tinker API Closed-Source Nature}
% ---------------------------------------------------------------------------

A central tension in \ours{} is that we benchmark \emph{against} a
closed-source platform while simultaneously advocating for reproducibility.
We include Tinker API experiments for three reasons: (i) Tinker represents
a realistic production environment for practitioners who lack on-premises
GPU clusters; (ii) the platform/open-source gap is itself a scientifically
interesting quantity to measure; and (iii) the high failure rate (11 of 14
experiments, 79\%) is a finding worth documenting publicly, as it reveals
the fragility of API-dependent training pipelines.

We make the following commitments to mitigate the reproducibility gap:
\begin{enumerate}
  \item All Tinker API experiment scripts, configuration files, and raw
        logs are archived in the repository.  Researchers with Tinker
        access can attempt replication using our exact configurations.
  \item All figures derived \emph{solely} from Tinker data are marked with
        ``\dag{}'' and accompanied by a note that independent reproduction
        requires Tinker API access.
  \item Primary statistical conclusions are drawn from Modal experiments
        wherever possible.  Tinker-only results are reported descriptively
        and are not used to support quantitative claims.
  \item If Tinker releases an open-source version or a compatible open
        alternative becomes available, we commit to re-running failed
        experiments and issuing a revised version of this paper.
\end{enumerate}

We believe honest disclosure of these constraints is more valuable to the
community than suppressing the Tinker experiments or presenting only
successful runs.
 into main.tex
% Intended placement: after the Conclusion section, before References.
% Label: \label{sec:impact}
% =============================================================================

\section{Broader Impact and Ethics Statement}
\label{sec:impact}

\subsection{Broader Impact}

\paragraph{Positive impacts.}
\ours{} lowers the barrier to reproducible RL post-training research by
providing a unified benchmark harness that works across multiple compute
platforms.  Prior work frequently reports results on a single proprietary
GPU cluster, making it difficult to determine how much of the observed
performance is attributable to the training algorithm versus the hardware
and scheduling environment.  By running identical experiments on both a
closed API (Tinker) and an open, explicitly provisioned cloud (Modal), we
surface \emph{platform-dependent variance} as a first-class experimental
variable.  This transparency allows practitioners to make more informed
decisions when selecting infrastructure for RL fine-tuning and helps the
community calibrate trust in benchmark results from single-platform papers.

Releasing all code, model checkpoints, and experiment logs publicly enables
independent replication, third-party auditing, and incremental improvement
without the need to re-run expensive GPU-hours from scratch.  We expect the
benchmark to be most valuable for researchers at resource-constrained
institutions who need a principled starting point for RL post-training
experiments.

\paragraph{Potential negative impacts and limitations.}
\begin{itemize}
  \item \textbf{Reward hacking.}  RLHF and GRPO optimise a proxy reward
        signal rather than the true underlying objective.  Models trained
        to maximise verifier-based rewards on GSM8K or HumanEval may learn
        to exploit idiosyncrasies of the verification function rather than
        genuinely improving reasoning or coding ability.  We include KL
        divergence tracking precisely to monitor policy drift, but reward
        hacking remains a risk in any production deployment of RL-fine-tuned
        models.

  \item \textbf{Alignment concerns.}  Reinforcement learning from human or
        proxy feedback can shift model behaviour in directions that are
        difficult to predict or reverse.  Models trained with poorly
        specified reward signals may become better at appearing correct while
        simultaneously becoming less calibrated or more prone to confident
        errors.  Our benchmark focuses on mathematically verifiable tasks,
        which partially mitigates this risk, but we caution against
        uncritical transfer of RL post-training pipelines to open-ended or
        subjective task domains without additional alignment safeguards.

  \item \textbf{Compute access disparities.}  Even with cost-optimised
        LoRA fine-tuning, reproducing the full \ours{} benchmark requires
        access to NVIDIA H100 GPUs (or equivalent) for 200--600 GPU-hours.
        This remains out of reach for many research groups, particularly in
        the Global South.  We partially address this by providing pretrained
        checkpoints, but the barrier to running novel experiments remains
        high.

  \item \textbf{Closed-platform dependency.}  A portion of our benchmark
        relies on the Tinker API, a proprietary, closed-source service.
        Results from Tinker experiments cannot be independently reproduced
        by researchers without API access.  We consider this a significant
        limitation and flag all Tinker-derived results accordingly in our
        figures.  We strongly recommend that readers treat Modal/open-source
        results as the primary reproducibility reference.
\end{itemize}

% ---------------------------------------------------------------------------
\subsection{Compute Costs and Carbon Footprint}
% ---------------------------------------------------------------------------

Table~\ref{tab:carbon} summarises our compute allocation and estimated
carbon footprint.  Carbon intensity estimates follow the methodology of
\citet{patterson2024empirical}, using a conservative average grid intensity
of 0.233 kg\,CO$_2$-eq per kWh for the United States (EPA eGRID 2023
average) and a power usage effectiveness (PUE) of 1.1 for cloud data
centres.  NVIDIA H100 SXM5 TDP is 700\,W; NVIDIA L4 TDP is 72\,W.

\begin{table}[h]
\centering
\caption{Compute allocation and estimated carbon footprint by platform.
  Tinker API carbon figures are marked N/A because GPU type is undisclosed.
  Modal figures use confirmed hardware specifications.}
\label{tab:carbon}
\begin{tabular}{llcccc}
\toprule
\textbf{Platform} & \textbf{GPU} & \textbf{Est. GPU-h} &
\textbf{TDP (W)} & \textbf{Est. kWh} &
\textbf{Est. CO\textsubscript{2} (kg)} \\
\midrule
Modal (H100 SXM5) & H100 SXM5 & $\sim$250 & 700 & $\sim$192 & $\sim$0.49 \\
Modal (L4, older runs) & L4       & $\sim$50  & 72  & $\sim$4   & $\sim$0.01 \\
Tinker API            & unknown  & $\sim$950 & N/A & N/A       & N/A        \\
\midrule
\multicolumn{3}{l}{\textbf{Total (Modal only, estimable)}}      & & $\sim$196 & $\sim$0.50 \\
\bottomrule
\end{tabular}
\end{table}

The estimated carbon footprint for the Modal portion of our experiments is
approximately \textbf{0.4--0.6\,kg\,CO$_2$-eq}, comparable to a short
car journey.  We cannot estimate the Tinker API footprint due to the
platform's opacity.  We acknowledge that the full experimental campaign
(including failed runs) consumed additional GPU-hours; our estimates above
reflect \emph{successful} jobs only.

% ---------------------------------------------------------------------------
\subsection{Data Provenance}
% ---------------------------------------------------------------------------

All training and evaluation datasets used in \ours{} are publicly available
research datasets.  No private, proprietary, or sensitive data were used
at any stage.

\begin{itemize}
  \item \textbf{GSM8K}~\cite{cobbe2021training}: 8,500 grade-school math
        word problems with natural-language solutions, released by OpenAI
        under the MIT licence.  We use the standard train/test splits
        without modification.

  \item \textbf{HumanEval}: 164 Python programming problems with unit
        tests, released by OpenAI under the MIT licence.  We use the
        original evaluation protocol (pass@$k$) without modification.

  \item \textbf{Synthetic tool-use dataset}: Generated by the authors from
        publicly documented APIs (calculator, web-search stub, calendar).
        No third-party proprietary data or user data were used in
        generation.  The generation prompts and all resulting examples are
        included in the repository under MIT licence.
\end{itemize}

No web-scraped data, personally identifiable information, or data subject
to copyright restrictions were used.  No human annotators were employed.

% ---------------------------------------------------------------------------
\subsection{Limitation Acknowledgment: Tinker API Closed-Source Nature}
% ---------------------------------------------------------------------------

A central tension in \ours{} is that we benchmark \emph{against} a
closed-source platform while simultaneously advocating for reproducibility.
We include Tinker API experiments for three reasons: (i) Tinker represents
a realistic production environment for practitioners who lack on-premises
GPU clusters; (ii) the platform/open-source gap is itself a scientifically
interesting quantity to measure; and (iii) the high failure rate (11 of 14
experiments, 79\%) is a finding worth documenting publicly, as it reveals
the fragility of API-dependent training pipelines.

We make the following commitments to mitigate the reproducibility gap:
\begin{enumerate}
  \item All Tinker API experiment scripts, configuration files, and raw
        logs are archived in the repository.  Researchers with Tinker
        access can attempt replication using our exact configurations.
  \item All figures derived \emph{solely} from Tinker data are marked with
        ``\dag{}'' and accompanied by a note that independent reproduction
        requires Tinker API access.
  \item Primary statistical conclusions are drawn from Modal experiments
        wherever possible.  Tinker-only results are reported descriptively
        and are not used to support quantitative claims.
  \item If Tinker releases an open-source version or a compatible open
        alternative becomes available, we commit to re-running failed
        experiments and issuing a revised version of this paper.
\end{enumerate}

We believe honest disclosure of these constraints is more valuable to the
community than suppressing the Tinker experiments or presenting only
successful runs.
